\begin{table*}[t]
\centering
\caption{Effect of additional positive-label information: budget 0.20 minus 0.05, in percentage points. Differences are computed between within-dataset split means before averaging equally across eight datasets.}
\label{tab:s-real-budget-effects}
\small
\setlength{\tabcolsep}{3.5pt}
\renewcommand{\arraystretch}{1.08}
\begin{tabularx}{\textwidth}{>{\raggedright\arraybackslash}Xrrrr}
\toprule
Method & AUROC OS & AUROC TS & AP OS & AP TS \\
\midrule
Logistic probability correction & +2.970 (8/8) & +2.745 (8/8) & +3.559 (8/8) & +3.241 (8/8) \\
nnPU logistic & +2.373 (8/8) & +2.220 (8/8) & +2.841 (8/8) & +2.454 (8/8) \\
Naive CatBoost & +2.473 (8/8) & +2.202 (8/8) & +3.627 (8/8) & +3.066 (8/8) \\
CS CatBoost & +3.015 (8/8) & +2.497 (8/8) & +4.456 (8/8) & +3.569 (8/8) \\
RN CatBoost & +1.691 (7/8) & +1.492 (8/8) & +1.903 (7/8) & +1.713 (7/8) \\
PU gradient boosting & +5.957 (8/8) & +3.232 (8/8) & +8.378 (8/8) & +4.719 (8/8) \\
\bottomrule
\end{tabularx}
\par\medskip
\begin{minipage}{\textwidth}\small
Parentheses give the number of datasets with a positive change. Each budget has 20 archived splits per condition; differences compare budget-specific means and are not assumed to pair identical observations across budgets. Dataset-specific values and negative/zero counts are retained in the CSV. No population-level inference is attached.
\end{minipage}
\end{table*}
